\documentclass[11pt,a4paper]{article}
\usepackage[utf8]{inputenc}
\usepackage[T1]{fontenc}
\usepackage{lmodern}
\usepackage[margin=1in]{geometry}
\usepackage{amsmath,amssymb,amsthm}
\usepackage{booktabs}
\usepackage{graphicx}
\usepackage{microtype}
\usepackage[numbers,sort&compress]{natbib}
\usepackage{xcolor}
\usepackage[colorlinks=true,linkcolor=blue!50!black,citecolor=blue!50!black,urlcolor=blue!50!black]{hyperref}
\usepackage{enumitem}
\usepackage{array}
\usepackage{caption}
\usepackage{float}
\usepackage{multicol}
\captionsetup{font=small,labelfont=bf}
\providecommand{\bibfont}{}
\renewcommand{\bibfont}{\footnotesize}
\setlength{\bibsep}{0.5pt plus 0.5pt}
\renewcommand{\floatpagefraction}{0.75}
\renewcommand{\topfraction}{0.9}
\renewcommand{\textfraction}{0.07}

\IfFileExists{numbers.tex}{% generated by experiments/make_numbers.py -- do not edit
% per-method E3 macros (Runs/Voronoi/Hartigan/Reach/MedianExcess) are generated systematically; main.tex uses a subset
\newcommand{\MetaSeed}{20260930}
\newcommand{\MetaSeedRelax}{20260931}
\newcommand{\MetaPython}{3.11.15}
\newcommand{\MetaNumpy}{2.4.6}
\newcommand{\MetaScipy}{1.17.1}
\newcommand{\MetaSklearn}{1.9.1}
\newcommand{\IdentitySeconds}{159}
\newcommand{\RelaxSeconds}{128}
\newcommand{\Nsynth}{300}
\newcommand{\Ndigits}{300}
\newcommand{\Niris}{150}
\newcommand{\NumDatasets}{5}
\newcommand{\NumPotentials}{6}
\newcommand{\EoneChecks}{19285}
\newcommand{\EonePassed}{19285}
\newcommand{\EoneMaxRelErr}{\ensuremath{1.1\times 10^{-11}}}
\newcommand{\EoneHookeSSEErr}{\ensuremath{3.9\times 10^{-16}}}
\newcommand{\EoneRandomPartitions}{200}
\newcommand{\EoneEqualSizeErr}{\ensuremath{2.2\times 10^{-16}}}
\newcommand{\EoneRatioMin}{29.7}
\newcommand{\EoneRatioMax}{197.5}
\newcommand{\EoneConfigs}{30}
\newcommand{\EonePartitionsTotal}{6090}
\newcommand{\EoneCndConfigs}{25}
\newcommand{\EoneCndPsd}{25}
\newcommand{\EoneRestConfigs}{5}
\newcommand{\EoneRestIndefinite}{5}
\newcommand{\EoneCndWorstRatio}{\ensuremath{3.2\times 10^{-16}}}
\newcommand{\EoneMaxSigma}{\ensuremath{4.4\times 10^{4}}}
\newcommand{\EtwoRuns}{90}
\newcommand{\EtwoIdentical}{90}
\newcommand{\EtwoMaxRelErr}{\ensuremath{1.4\times 10^{-14}}}
\newcommand{\EtwoNsub}{120}
\newcommand{\EtwoAnchoredRuns}{15}
\newcommand{\EtwoAnchoredIdentical}{15}
\newcommand{\EtwoAnchoredMaxErr}{\ensuremath{4.3\times 10^{-16}}}
\newcommand{\EtwoSecondsPhysics}{56}
\newcommand{\EtwoSecondsKernel}{0.6}
\newcommand{\EtwoMaxSweeps}{11}
\newcommand{\EtwoMeanSweeps}{4.9}
\newcommand{\EfourN}{7}
\newcommand{\EfourK}{2}
\newcommand{\EfourPartitions}{63}
\newcommand{\EfourUnknowns}{22}
\newcommand{\EfourSeed}{20260932}
\newcommand{\EfourPoints}{(5,10), (6,2), (4,4), (2,2), (3,7), (0,6), (8,8)}
\newcommand{\EfourPointsB}{(0,2), (7,4), (8,7), (5,0), (2,8), (2,3), (4,10)}
\newcommand{\EfourCells}{12}
\newcommand{\EfourCellsSame}{12}
\newcommand{\EfourRankSpecific}{21}
\newcommand{\EfourRankTotal}{22}
\newcommand{\EthreeConfigs}{30}
\newcommand{\EthreeR}{10}
\newcommand{\EthreeRanneal}{5}
\newcommand{\EthreeAnnealSweeps}{40}
\newcommand{\EthreeRunsRelax}{600}
\newcommand{\EthreeVoronoiRelax}{600}
\newcommand{\EthreeHartiganRelax}{600}
\newcommand{\EthreeReachRelax}{28}
\newcommand{\EthreeMedianExcessRelax}{\ensuremath{8.4\times 10^{-4}}}
\newcommand{\EthreeRunsLloydRelax}{600}
\newcommand{\EthreeVoronoiLloydRelax}{600}
\newcommand{\EthreeHartiganLloydRelax}{600}
\newcommand{\EthreeReachLloydRelax}{25}
\newcommand{\EthreeMedianExcessLloydRelax}{\ensuremath{7.7\times 10^{-4}}}
\newcommand{\EthreeRunsLloyd}{600}
\newcommand{\EthreeVoronoiLloyd}{600}
\newcommand{\EthreeHartiganLloyd}{204}
\newcommand{\EthreeReachLloyd}{14}
\newcommand{\EthreeMedianExcessLloyd}{\ensuremath{3.7\times 10^{-3}}}
\newcommand{\EthreeRunsAnneal}{150}
\newcommand{\EthreeVoronoiAnneal}{150}
\newcommand{\EthreeHartiganAnneal}{150}
\newcommand{\EthreeReachAnneal}{20}
\newcommand{\EthreeMedianExcessAnneal}{\ensuremath{1.7\times 10^{-12}}}
\newcommand{\EthreeRunsNaive}{600}
\newcommand{\EthreeVoronoiNaive}{600}
\newcommand{\EthreeHartiganNaive}{38}
\newcommand{\EthreeReachNaive}{6}
\newcommand{\EthreeMedianExcessNaive}{\ensuremath{5.9\times 10^{-1}}}
\newcommand{\EthreeDistinguishable}{6}
\newcommand{\EthreeRelaxBetter}{5}
\newcommand{\EthreeLloydRelaxBetter}{1}
\newcommand{\EthreeMaxAbsDiff}{\ensuremath{7.4\times 10^{-3}}}
\newcommand{\EthreeAriOne}{24}
\newcommand{\EthreeDirectCheck}{\ensuremath{4.3\times 10^{-13}}}
\newcommand{\EthreeShiftedConfigs}{5}
\newcommand{\EthreeLloydHartiganPct}{34}
\newcommand{\EthreeNaiveHartiganPct}{6}
\newcommand{\EthreeAgreeButDifferent}{0}
}{}

\theoremstyle{plain}
\newtheorem{theorem}{Theorem}[section]
\newtheorem{proposition}[theorem]{Proposition}
\newtheorem{lemma}[theorem]{Lemma}
\newtheorem{conjecture}[theorem]{Conjecture}
\theoremstyle{definition}
\newtheorem{definition}[theorem]{Definition}
\newtheorem{problem}[theorem]{Open problem}
\theoremstyle{remark}
\newtheorem{remark}[theorem]{Remark}

\newcommand{\R}{\mathbb{R}}
\newcommand{\Q}{\mathbb{Q}}
\newcommand{\dd}{\mathrm{d}}
\newcommand{\tr}{\operatorname{tr}}
\newcommand{\SSE}{\mathrm{SSE}}
\newcommand{\spec}{\mathrm{sp}}
\newcommand{\tot}{\mathrm{tot}}
\newcommand{\one}{\mathbf 1}
\newcommand{\status}[1]{\textsf{\small [#1]}}

\title{\textbf{Spring-Energy Clustering Is Kernel Clustering:}\\[2pt]
\large Known Equivalences Behind a Null Screening}
\author{Luciano Silva Alarco\\
\small Pontificia Universidad Cat\'olica del Per\'u\\
\small ORCID 0000-0002-3395-3129}
\date{\small Working draft v0.3 (two rounds of internal review applied) --- research line LMI001 --- 3 October 2026}

\begin{document}
\maketitle

\begin{abstract}
\noindent
A research line on ``mechanical analogies for data clustering'' screened objectives built from springs, tension and lifting against matched kernel controls and found no surviving mechanical component. No manuscript or code of that screening was available to us: we reconstruct the family from the line's public summary as \emph{pairwise} spring energies (a pair potential $\phi(d)$, optionally after a lift, normalised per particle or summed) and read the matched control as kernel $k$-means on the same matrix. For that family the null result follows from known results, collected here in the language of springs: the per-particle energy is the pairwise-clustering cost of Hofmann and Buhmann with dissimilarities $\phi(d_{ij})$, hence, up to a constant and a factor $\tfrac12$, a kernel $k$-means objective that a constant shift makes positive semidefinite without changing the minimisers (Roth et al.; Dhillon, Guan and Kulis); Hooke springs give $k$-means; for potentials of negative type, such as the constant-tension string $\phi(d)=d$, the energy is the within-cluster dispersion of energy statistics and the kernel is the dataset-independent distance-induced kernel of Sejdinovic et al.\ (Fran\c{c}a, Rizzo and Vogelstein); the total energy and its tension dual are min-sum $k$-clustering and max-$k$-cut; and zero-temperature relaxation is Hartigan's method in kernel space (kernel $k$-groups), whose fixed points are Voronoi-stable, i.e.\ Lloyd fixed points with ties broken in favour of the current cluster. Our own contribution is small: the synthesis and its reading for the line; the explicit mechanism by which a diagonal shift freezes Lloyd's rule; a proof that the rest-length spring is not of negative type; and exact rational certificates, on two explicit configurations, that a triangle-area three-body term, a cluster-mean rest length and a per-spring average are not pairwise objectives under the per-particle or total normalisation (a statement about objective values, not minimisers). Computationally, on five datasets and six potentials the spring energies and the kernel objectives coincide to $\EoneMaxRelErr$ on \EonePartitionsTotal{} partitions, spring and kernel implementations of the same relaxation follow identical trajectories in \EtwoIdentical/\EtwoRuns{} runs, and all \EthreeRunsRelax{} relaxation fixed points are kernel $k$-means fixed points. Dynamics on positions is left open. No claim is made about clustering performance.
\end{abstract}

\section{Introduction}\label{sec:intro}

Physical metaphors recur in clustering: algorithms have been built on gravity~\citep{wright1977}, on spin models at finite temperature~\citep{blatt1996}, on deterministic annealing of a free energy~\citep{rose1990,rose1998}, on quantum potentials~\citep{horn2002}, and on elastic nets~\citep{durbin1987,yuille1990}. The research line to which this note belongs (LMI001, ``Mechanical analogies for data clustering'') studied formulations of a more elementary kind: points connected by springs, with tension and lifting terms, clustered by minimising an elastic energy. Its stated objective was to determine whether the mechanical component produces an effect distinguishable from matched kernel controls, and its screening found none. \emph{The definitions of Section~\ref{sec:family} are a reconstruction from the line's public summary; no manuscript or code of the original screening was available to us, and the conclusions below apply to objectives of that pairwise form.} We likewise read ``matched kernel control'' as kernel $k$-means on the same affinity matrix with the same optimiser; this too is a reconstruction. The purpose of the note is to explain the null result rather than to repeat it, and the explanation is not new: for objectives of the pairwise form it is already contained in three literatures that the line's summary does not mention. Pairwise data clustering~\citep{hofmann1997} minimises exactly the per-particle spring energy for a dissimilarity matrix (by deterministic annealing, a finite-temperature treatment of the same energy); its reduction to $k$-means after a constant shift is the constant-shift embedding of \citet{roth2003}; the same objective is the ratio association that weighted kernel $k$-means optimises~\citep{dhillon2005tr,dhillon2007}. For potentials of negative type, energy statistics supplies a dataset-independent kernel~\citep{sejdinovic2013}, and clustering by within-cluster energy is kernel $k$-means, optimised preferably by Hartigan's method~\citep{franca2021}. A control that matches the kernel therefore matches the objective exactly, and no experiment on objective values can separate them (experiments on optimisers can, Section~\ref{sec:experiments}).

\paragraph{Contribution.} The note collects these results for springs, with explicit constants, and draws the consequence for the research line (Section~\ref{sec:results}). Its own additions are modest: the explicit mechanism by which a diagonal shift freezes Lloyd's rule while the energy and the single-move relaxation are unchanged (Proposition~\ref{prop:shift}), quantified in E3; a short proof that the rest-length spring is not of negative type (Proposition~\ref{prop:bernstein}); exact rational certificates that three specific functionals (a triangle-area three-body term, a cluster-mean rest length and a per-spring average) are not pairwise objectives under the per-particle or total normalisation (Section~\ref{sec:escape}; these concern objective values, not minimisers); and a computational confirmation of the identities with predefined criteria (Section~\ref{sec:experiments}). The new mathematical content is small; the note's value, if any, is to make the null screening a corollary of known results.

\paragraph{Scope.} The note is about objectives and their minimisers, not about clustering quality. We do not claim that any potential clusters real data well, and the adjusted Rand indices reported are context, not results. The datasets are standard synthetic ones and two small scikit-learn datasets~\citep{pedregosa2011}. Every number in the text is generated by the scripts described in Appendix~\ref{app:repro}, with fixed seeds. Section~\ref{sec:claims} lists all claims with their status.

\section{The pairwise spring family}\label{sec:family}

Throughout, $X=(x_1,\dots,x_n)$ is a configuration in $\R^d$, $k\ge2$ is fixed, and a \emph{partition} $C=(C_1,\dots,C_k)$ is a partition of $[n]$ into $k$ non-empty clusters, with sizes $n_c=|C_c|$ and indicator vectors $\one_c\in\{0,1\}^n$. We write $D_{ij}=\lVert x_i-x_j\rVert$. ``PSD'' always refers to a matrix (positive semidefinite); a \emph{kernel} on a set is called \emph{positive definite} (PD) when all its Gram matrices are PSD, the convention of \citet{bcr1984}.

\begin{definition}[Spring network, elastic energy, variants]\label{def:family}
A \emph{spring network} is a triple $(\phi,L,w)$ of a \emph{pair potential} $\phi\colon[0,\infty)\to\R$, a \emph{lift} $L\colon\R^d\to\R^D$ (the identity by default), and a \emph{normalisation} $w\in\{w_\spec,w_\tot\}$ with $w_\spec(m)=1/m$ (energy per particle of the cluster) and $w_\tot(m)=1$ (total energy). Its \emph{within-cluster elastic energy} is
\begin{equation}\label{eq:E}
E_{\phi,L,w}(C)\;=\;\sum_{c=1}^k w(n_c)\sum_{\substack{i<j\\ i,j\in C_c}}\phi\bigl(\lVert L x_i-L x_j\rVert\bigr).
\end{equation}
We write $A_{ij}=\phi(\lVert Lx_i-Lx_j\rVert)$ for the \emph{spring matrix}, $E_\spec$ and $E_\tot$ for the two normalisations, and $\SSE(C)=\sum_c\sum_{i\in C_c}\lVert x_i-\mu_c\rVert^2$, $\mu_c$ the centroid, for the $k$-means objective~\citep{macqueen1967,lloyd1982}. The \emph{released tension} of $C$ is the energy of the springs that cross clusters, $T(C)=\sum_{c<c'}\sum_{i\in C_c,\,j\in C_{c'}}A_{ij}$; a tension formulation maximises $T$. An \emph{anchored} network attaches every particle to a hub $y_c\in\R^D$ of its cluster, with energy $E^{\mathrm{anch}}(C,y)=\sum_c\sum_{i\in C_c}\phi(\lVert Lx_i-y_c\rVert)$, minimised over the hubs as well as over $C$.
\end{definition}

The potentials used below are: the Hooke spring with zero rest length, $\phi(d)=d^2$ (the spring constant is absorbed); the Hooke spring with rest length $r>0$, $\phi(d)=(d-r)^2$; the string under constant tension, $\phi(d)=d$, whose per-particle energy is the within-cluster dispersion of energy statistics~\citep{franca2021}; the saturating (breakable) spring, $\phi(d)=1-e^{-d^2/2\ell^2}$; the soft spring, $\phi(d)=\log(1+d^2/\ell^2)$; and the Hooke spring on the paraboloid lift $L(x)=(x,\lVert x\rVert^2/s)$.

\begin{definition}[Kernel $k$-means objective, trace form]\label{def:kkm}
For a symmetric matrix $K\in\R^{n\times n}$ and a partition $C$,
\begin{equation}\label{eq:J}
J_K(C)\;=\;\tr K-\sum_{c=1}^k\frac{\one_c^{\!\top}K\one_c}{n_c}.
\end{equation}
When $K$ is PSD with a Gram factorisation $K_{ij}=\langle\Phi_i,\Phi_j\rangle$, $J_K(C)=\sum_c\sum_{i\in C_c}\lVert\Phi_i-m_c\rVert^2$ with $m_c$ the centroid of the $\Phi_i$, $i\in C_c$: the $k$-means objective in feature space~\citep{scholkopf1998,girolami2002,zha2001,dhillon2004}. For general symmetric $K$ we take~\eqref{eq:J} as the definition.
\end{definition}

\begin{definition}[Negative type, Bernstein functions]\label{def:cnd}
A symmetric kernel $\psi$ on a set $S$ is \emph{conditionally negative definite} (CND) if $\sum_{i,j}c_ic_j\psi(s_i,s_j)\le0$ for all finite families $s_i\in S$ and reals $c_i$ with $\sum_ic_i=0$. A function $\phi\colon[0,\infty)\to\R$ is \emph{of negative type} if $(x,y)\mapsto\phi(\lVert x-y\rVert)$ is CND on $\R^d$ for every $d$. A function $\psi\colon[0,\infty)\to[0,\infty)$ is a \emph{Bernstein function} if it is continuous, $C^\infty$ on $(0,\infty)$, and $(-1)^{m-1}\psi^{(m)}\ge0$ for all $m\ge1$; equivalently (L\'evy--Khintchine representation, \citealp[Thm.~3.2]{ssv2012}) $\psi(s)=a+bs+\int_{(0,\infty)}(1-e^{-ts})\,\mu(\dd t)$ with $a,b\ge0$ and $\int\min(1,t)\,\mu(\dd t)<\infty$.
\end{definition}

\section{Results}\label{sec:results}

\subsection{Invariances of the trace form}

\begin{lemma}[What $J$ does not see]\label{lem:inv}
Let $K$ be symmetric, $C$ a partition into $k$ non-empty clusters, $\sigma,a\in\R$ and $f\in\R^n$. Then
\[
J_{K+\sigma I}(C)=J_K(C)+\sigma(n-k),\qquad J_{K+a\one\one^{\!\top}}(C)=J_K(C),\qquad J_{K+f\one^{\!\top}+\one f^{\!\top}}(C)=J_K(C).
\]
\end{lemma}
\begin{proof}
$\tr(\sigma I)=\sigma n$ and $\one_c^{\!\top}\sigma I\one_c/n_c=\sigma$ for each of the $k$ clusters. $\tr(a\one\one^{\!\top})=an$ and $\sum_c a n_c^2/n_c=an$. $\tr(f\one^{\!\top}+\one f^{\!\top})=2\sum_if_i$ and $\one_c^{\!\top}(f\one^{\!\top}+\one f^{\!\top})\one_c/n_c=2\sum_{i\in C_c}f_i$, which sums to $2\sum_i f_i$ over $c$.
\end{proof}

In particular the minimisers of $J_K$ over partitions into $k$ clusters are unchanged by a diagonal shift, by a constant added to every entry, and by a separable term $f_i+f_j$. The lemma is elementary and not new: the diagonal shift is the constant-shift embedding of \citet{roth2003} (a constant added to the off-diagonal dissimilarities) and the device of \citet{dhillon2005tr,dhillon2007} for making an indefinite affinity matrix PSD without changing the optimal partitions; the other two invariances connect the spring energy to Schoenberg's theory of metric embeddings~\citep{schoenberg1935,schoenberg1938} below.

\subsection{Springs are kernels}

\begin{theorem}[Spring energies are kernel objectives; collected from the literature]\label{thm:main}
Let $(\phi,L,w)$ be a spring network with spring matrix $A$ and let $K=-A$ (diagonal included, $K_{ii}=-\phi(0)$). For every partition $C$ into $k$ non-empty clusters:
\begin{enumerate}[label=(\alph*),leftmargin=2em]
\item $E_{\phi,L,w}(C)$ depends on the data only through $A$.
\item (per-particle normalisation) $E_\spec(C)=\tfrac12 J_K(C)+\tfrac12(n-k)\phi(0)$. For every $\sigma\ge-\lambda_{\min}(K)$ the matrix $K+\sigma I$ is PSD and $E_\spec(C)=\tfrac12J_{K+\sigma I}(C)+\tfrac12(n-k)(\phi(0)-\sigma)$; hence $\arg\min_C E_\spec=\arg\min_C J_{K+\sigma I}$, the $k$-means problem in the feature space of any Gram factorisation of $K+\sigma I$.
\item (Hooke) If $\phi(d)=d^2$ and $L$ is the identity, $E_\spec(C)=\SSE(C)$ exactly.
\item (total normalisation and tension) $E_\tot(C)+T(C)=\sum_{i<j}A_{ij}$ is independent of $C$; minimising $E_\tot$ is maximising the released tension, which is the max-$k$-cut, equivalently min-sum $k$-clustering, problem with weights $A_{ij}$~\citep{sahni1976}. If all clusters have the same size, $E_\tot=(n/k)E_\spec$.
\item (converse) Every kernel $k$-means objective is a Hooke energy: if $K$ is PSD and $K=\Phi\Phi^{\!\top}$, then $J_K(C)=E_{d^2,\Phi,\spec}(C)$, the per-particle Hooke energy of the lifted configuration $(\Phi_1,\dots,\Phi_n)$.
\item (universal kernel) If $\phi$ is of negative type, then
\[
\widetilde K(x,y)=\phi(\lVert x\rVert)+\phi(\lVert y\rVert)-\phi(\lVert x-y\rVert)-\phi(0)
\]
is a PD kernel on $\R^D$ that does not depend on the dataset (given the scales $\ell,r,s$ that enter $\phi$ and $L$), and $E_\spec(C)=\tfrac12J_{\widetilde K}(C)+\tfrac12(n-k)\phi(0)$ where $\widetilde K$ is evaluated on the lifted points. If $\Phi$ is a feature map of $\widetilde K$ and $\widetilde\Phi=\Phi/\sqrt2$, then $\lVert\widetilde\Phi(x)-\widetilde\Phi(y)\rVert^2=\phi(\lVert x-y\rVert)-\phi(0)$ and
\[
E_\spec(C)=E_{d^2,\widetilde\Phi\circ L,\spec}(C)+\tfrac12(n-k)\phi(0):
\]
up to the constant, the spring energy is the per-particle Hooke energy of the lifted configuration $(\widetilde\Phi(Lx_i))_i$ (the lift $\widetilde\Phi\circ L$ takes values in a possibly infinite-dimensional feature space; restricted to the $n$ lifted points it has finite rank, so Definition~\ref{def:family} applies with finite $D$). Without the rescaling, $E_\spec$ is \emph{half} the Hooke energy of $(\Phi(Lx_i))_i$ plus the constant.
\end{enumerate}
\end{theorem}

\noindent\emph{Status.} The theorem collects known results; we state it for springs because the explicit constants are what E1 tests. (a), (d) elementary; (c) classical; (e) the feature-space reading of $k$-means~\citep{zha2001,girolami2002} \status{classical}. (b) \status{literature}: with zero diagonal, $E_\spec$ is the pairwise-clustering cost of \citet{hofmann1997} for the dissimilarities $A_{ij}$, and its reduction to $k$-means after a constant shift is the constant-shift embedding of \citet{roth2003}; equivalently, $\sum_c\one_c^{\!\top}A\one_c/n_c$ is the ratio association of $A$, whose maximisation is kernel $k$-means with $K=\sigma I+A$~\citep{dhillon2005tr,dhillon2007}, and minimising $E_\spec$ maximises the ratio association of $-A$, i.e.\ uses $K=\sigma I-A$. Specific to springs is only the constant $\tfrac12(n-k)\phi(0)$. (f) \status{literature}: $\tfrac12\widetilde K$ is the distance-induced kernel of \citet{sejdinovic2013} for the semimetric of negative type $\rho(x,y)=\phi(\lVert x-y\rVert)-\phi(0)$ with base point $0$ (it is $\ge0$ because the CND test on two points forces $\phi(d)\ge\phi(0)$), $\widetilde\Phi$ is its feature map, $E_\spec-\tfrac12(n-k)\phi(0)$ is the within-cluster energy dispersion of energy statistics for $\rho$, and its equivalence with kernel $k$-means is in \citet{franca2021}; we include the short argument.

\begin{proof}
(a) is the definition. (b) Since $\sum_{i<j\in C_c}A_{ij}=\tfrac12(\one_c^{\!\top}A\one_c-n_c\phi(0))$,
\[
E_\spec(C)=\tfrac12\sum_c\frac{\one_c^{\!\top}A\one_c}{n_c}-\tfrac12 k\phi(0),
\qquad\text{and}\qquad
\sum_c\frac{\one_c^{\!\top}A\one_c}{n_c}=-\sum_c\frac{\one_c^{\!\top}K\one_c}{n_c}=J_K(C)-\tr K,
\]
with $\tr K=-n\phi(0)$, so $E_\spec(C)=\tfrac12J_K(C)+\tfrac12(n-k)\phi(0)$. The shifted form follows from Lemma~\ref{lem:inv}, and $K+\sigma I\succeq0$ for $\sigma\ge-\lambda_{\min}(K)$. (c) For any finite set $P$ with centroid $\mu$, $\sum_{i,j\in P}\lVert x_i-x_j\rVert^2=2|P|\sum_{i\in P}\lVert x_i-\mu\rVert^2$ (expand $\lVert(x_i-\mu)-(x_j-\mu)\rVert^2$; the cross terms vanish). Halving and dividing by $|P|$ gives $\frac1{n_c}\sum_{i<j\in C_c}d_{ij}^2=\sum_{i\in C_c}\lVert x_i-\mu_c\rVert^2$; sum over $c$. (d) Every pair $\{i,j\}$ is either within a cluster or across, so $E_\tot+T$ is the total spring energy; the equal-size statement is $w_\tot=(n/k)w_\spec$. (e) By (b) with $\phi(0)=0$ applied to the configuration $\Phi$: $E_{d^2,\Phi,\spec}(C)=\tfrac12J_{-D_\Phi^{(2)}}(C)$, where $(D_\Phi^{(2)})_{ij}=K_{ii}+K_{jj}-2K_{ij}$, so $-D_\Phi^{(2)}=2K-f\one^{\!\top}-\one f^{\!\top}$ with $f_i=K_{ii}$, and Lemma~\ref{lem:inv} gives $\tfrac12J_{2K}(C)=J_K(C)$. (f) Positive definiteness of $\widetilde K$ is the lemma relating CND and PD kernels~\citep[Ch.~3, \S2, Lemma~2.1]{bcr1984}: for points $x_1,\dots,x_m$ and reals $c_i$, apply the CND inequality to the family augmented by $x_0=0$, $c_0=-\sum_{i\ge1}c_i$; substituting $c_0$ it reads $\sum_{i,j\ge1}c_ic_j\bigl[\phi(\lVert x_i-x_j\rVert)-\phi(\lVert x_i\rVert)-\phi(\lVert x_j\rVert)+\phi(0)\bigr]\le0$, i.e.\ $\sum_{i,j}c_ic_j\widetilde K(x_i,x_j)\ge0$. On the dataset, $\widetilde K=K+f\one^{\!\top}+\one f^{\!\top}-\phi(0)\one\one^{\!\top}$ with $f_i=\phi(\lVert Lx_i\rVert)$, so $J_{\widetilde K}=J_K$ by Lemma~\ref{lem:inv} and the identity follows from (b). Next, $\lVert\Phi(x)-\Phi(y)\rVert^2=\widetilde K(x,x)+\widetilde K(y,y)-2\widetilde K(x,y)=2\phi(\lVert x-y\rVert)-2\phi(0)$, so $\lVert\widetilde\Phi(x)-\widetilde\Phi(y)\rVert^2=\phi(\lVert x-y\rVert)-\phi(0)$. Finally (e) applied to $\widetilde K$ gives $J_{\widetilde K}(C)=E_{d^2,\Phi\circ L,\spec}(C)$, and scaling a configuration by $1/\sqrt2$ halves every squared distance, so $\tfrac12J_{\widetilde K}(C)=E_{d^2,\widetilde\Phi\circ L,\spec}(C)$. For Hooke springs $\widetilde K(x,y)=2\langle x,y\rangle$, $\Phi(x)=\sqrt2\,x$ and $\widetilde\Phi(x)=x$.
\end{proof}

\begin{proposition}[Which springs are of negative type]\label{prop:bernstein}
\begin{enumerate}[label=(\roman*),leftmargin=2em]
\item $\phi(d)=d^2$ is of negative type: for $\sum_ic_i=0$, $\sum_{i,j}c_ic_j\lVert x_i-x_j\rVert^2=-2\bigl\lVert\sum_ic_ix_i\bigr\rVert^2\le0$. \status{classical}
\item If $N$ is a CND kernel with $N\ge0$ and $\psi$ is a Bernstein function, then $\psi\circ N$ is CND. \status{literature}
\item Consequently $\phi(d)=\psi(d^2)$ is of negative type for every Bernstein $\psi$; this covers $d^{2\alpha}$ for $0<\alpha\le1$ (in particular $d$), $\log(1+d^2/\ell^2)$, $1-e^{-d^2/2\ell^2}$ and $d^2/(1+d^2/\ell^2)$. The rest-length potential $(d-r)^2$ is \emph{not} of negative type. \status{proved here (elementary)}
\end{enumerate}
\end{proposition}
\begin{proof}[Proof sketch]
(i) Expand $\lVert x_i-x_j\rVert^2=\lVert x_i\rVert^2+\lVert x_j\rVert^2-2\langle x_i,x_j\rangle$; the first two terms vanish against $\sum c_i=0$. (ii) is proved in \citet[Ch.~3, \S2]{bcr1984} (Schoenberg's theorem is Theorem~2.2 there; the composition with Bernstein functions follows from it and the representation in Definition~\ref{def:cnd}), see also \citet{ssv2012}; we do not vouch for the exact numbering of the composition result. The argument is: by Schoenberg's theorem~\citep{schoenberg1938} $N$ is CND iff $e^{-tN}$ is PD for all $t>0$, so $1-e^{-tN}$ is CND (a constant kernel is CND-neutral and the negative of a PD kernel is CND); write $\psi(N)=a+bN+\int(1-e^{-tN})\mu(\dd t)$ and note that non-negative combinations and pointwise limits of CND kernels are CND. (iii) The listed $\psi$ are Bernstein: $s^\alpha$ ($0<\alpha\le1$), $\log(1+s/\ell^2)$, $1-e^{-s/2\ell^2}$, $s/(1+s/\ell^2)$ all have completely monotone derivatives. If $\phi$ is of negative type, the CND inequality for two points with $c=(1,-1)$ gives $2\phi(0)-2\phi(d)\le0$, i.e.\ $\phi(d)\ge\phi(0)$ for all $d$; for $(d-r)^2$, $\phi(r)=0<r^2=\phi(0)$.
\emph{Rest length.} Theorem~\ref{thm:main}(b) still applies to $(d-r)^2$: its spring energy is the kernel $k$-means objective of $-(D-r)^{(2)}$ after a diagonal shift, and its minimisers are feature-space $k$-means minimisers for that dataset. What fails is a dataset-independent PD kernel. $\widetilde K$ is not PD, since $\widetilde K(x,x)=2(\phi(\lVert x\rVert)-\phi(0))=-2r^2<0$ when $\lVert x\rVert=r$. Nor is any kernel $\widetilde K(x,y)+f(x)+f(y)+a+\sigma[x=y]$ with $f$, $a$, $\sigma$ fixed: for the collinear points $0,su,2su,3su$ ($u$ a unit vector, $s>0$) and $c=(1,-1,-1,1)$, the terms in $f$ and $a$ vanish against $\sum_ic_i=0$ and the quadratic form of the Gram matrix is $\sum_{i,j}c_ic_j(-\phi(d_{ij}))+4\sigma=-8rs+4\sigma<0$ once $s>\sigma/2r$. Whether Lemma~\ref{lem:inv} lists all the transformations that leave $J$ unchanged up to a constant on every dataset, and hence whether some other dataset-independent PD representative exists, we leave open. Accordingly, E1 finds $\widetilde K$ indefinite for this potential on every dataset (Section~\ref{sec:experiments}).
\end{proof}

\begin{remark}[Lifting and anchoring]\label{rem:lift}
(a) $E_{\phi,L,w}(X)=E_{\phi,\mathrm{id},w}(LX)$, so a lift replaces $A$ by $\phi(\lVert Lx_i-Lx_j\rVert)$ and nothing else; for Hooke springs a lift is exactly the choice of the kernel $\langle Lx,Ly\rangle$, and by Theorem~\ref{thm:main}(e) every PD kernel arises this way on finite data \status{trivial}. (b) For Hooke springs, $\min_y E^{\mathrm{anch}}(C,y)=E_\spec(C)=\SSE(C)$, attained at the centroids (Theorem~\ref{thm:main}(c) and the fact that the centroid minimises the sum of squared distances), and the alternating equilibration ``hubs move to the equilibrium of their springs; particles re-attach to the hub of least energy'' is Lloyd's algorithm~\citep{lloyd1982} by definition. With $\phi(d)=d$ the anchored objective is the $k$-median objective (geometric medians as hubs) \status{classical}; we make no claim about whether such anchored objectives are of the form~\eqref{eq:E}.
\end{remark}

\subsection{Relaxation of the spring system}

The mechanical way to minimise~\eqref{eq:E} is to let the system relax: one particle at a time jumps to the cluster in which its springs store less energy, or, at finite temperature, jumps with Metropolis probability. On assignments this is a known algorithm.

\begin{proposition}[Zero-temperature relaxation is Hartigan's method]\label{prop:hartigan}
Let $K'$ be any PSD matrix with $J_{K'}=J_K+\mathrm{const}$ on partitions (for instance $K+\sigma I$ or, for $\phi$ of negative type, $\widetilde K$), and write $\delta_{K'}(i,c)=K'_{ii}-\tfrac2{n_c}(K'\one_c)_i+\tfrac1{n_c^2}\one_c^{\!\top}K'\one_c$ for the squared feature distance from $\Phi_i$ to the centroid of cluster $c$. Moving particle $i$ from its cluster $a$ ($n_a\ge2$) to cluster $c$ changes $E_\spec$ by
\begin{equation}\label{eq:delta}
\Delta E=\tfrac12\Bigl[\tfrac{n_c}{n_c+1}\,\delta_{K'}(i,c)-\tfrac{n_a}{n_a-1}\,\delta_{K'}(i,a)\Bigr],
\end{equation}
and the right-hand side does not depend on the representative $K'$. Consequently:
\begin{enumerate}[label=(\roman*),leftmargin=2em]
\item the greedy single-particle relaxation of the spring energy is Hartigan's method~\citep{hartigan1975,hartiganwong1979} applied to the kernel $k$-means objective;
\item every partition that is stable under single-particle moves is Voronoi-stable in the feature space of every PSD representative $K'$, i.e.\ each $\Phi_i$ is assigned to a nearest centroid, so it is a fixed point of Lloyd's algorithm for $K'$ with ties broken in favour of the current cluster (on the relation between such fixed points and local optimality see \citealp{selim1984});
\item the converse fails: $X=\{0,2,3.4,3.6\}\subset\R$, $k=2$, $C=\{\{0,2\},\{3.4,3.6\}\}$ is Voronoi-stable with $\SSE=2.02$, while moving the point $2$ gives $\SSE=1.52$.
\end{enumerate}
\end{proposition}
\begin{proof}
Write $S_{ic}=(K'\one_c)_i$ and $Q_c=\one_c^{\!\top}K'\one_c$. Removing $i$ from $a$ changes $Q_a/n_a$ to $(Q_a-2S_{ia}+K'_{ii})/(n_a-1)$, adding it to $c$ changes $Q_c/n_c$ to $(Q_c+2S_{ic}+K'_{ii})/(n_c+1)$, and $\Delta J=-\Delta\sum_cQ_c/n_c$. Expanding both differences and comparing with $\delta_{K'}$ gives $\Delta J=\frac{n_c}{n_c+1}\delta_{K'}(i,c)-\frac{n_a}{n_a-1}\delta_{K'}(i,a)$, an algebraic identity valid for any symmetric $K'$; $\Delta E=\tfrac12\Delta J$ by Theorem~\ref{thm:main}(b). Independence of the representative: if $J_{K'}=J_K+\mathrm{const}$ on partitions then $\Delta J_{K'}=\Delta J_K$ for every move, and the identity holds for both matrices, so the right-hand side of~\eqref{eq:delta} is the same for every such $K'$. Explicitly, under $K'\mapsto K'+\sigma I$, $\delta(i,c)$ increases by $\sigma(1+1/n_c)$ for $i\notin C_c$ and by $\sigma(1-1/n_a)$ for $i\in C_a$, and both terms of~\eqref{eq:delta} increase by $\sigma$; under the other two invariances of Lemma~\ref{lem:inv}, $\delta$ itself is unchanged. (i) is the definition of Hartigan's move. (ii) If no move decreases $E$ then $\frac{n_c}{n_c+1}\delta(i,c)\ge\frac{n_a}{n_a-1}\delta(i,a)$ for all $c\ne a$, and since $\frac{n_c}{n_c+1}<1<\frac{n_a}{n_a-1}$ this forces $\delta(i,c)\ge\delta(i,a)$ (not necessarily strictly, whence the tie rule). (iii) Centroids $1$ and $3.5$; the point $2$ is at distance $1$ from the first and $1.5$ from the second; after the move the clusters $\{0\}$ and $\{2,3.4,3.6\}$ (centroid $3$) have $\SSE=0+1+0.16+0.36=1.52$; by~\eqref{eq:delta} with $K'=\widetilde K=2\langle x,y\rangle$, $\Delta E=\tfrac12\bigl(\tfrac23\cdot2\cdot1.5^2-2\cdot2\cdot1^2\bigr)=-0.5$ (equivalently, the bracket of~\eqref{eq:delta}, which is $\Delta J_{K'}$ for any symmetric $K'$, with $K'=\langle x,y\rangle$ and $J_{\langle x,y\rangle}=\SSE$).
\end{proof}
\noindent \emph{Status.} (i) \status{literature}: Hartigan's method in kernel space is the kernel $k$-groups algorithm of \citet{franca2021}, there derived for energy statistics; the representative-free form~\eqref{eq:delta} is elementary. (ii) is the kernel form of the Euclidean inclusion of \citet{telgarsky2010} \status{proved here}; we could not check whether \citet{franca2021} state it and claim no priority. (iii) is the standard type of example. The consequence for the research line is that ``spring relaxation'' cannot converge to anything outside the set of kernel $k$-means fixed points.

\begin{proposition}[Lloyd's rule is representation-dependent]\label{prop:shift}
For $K'=K+\sigma I$, $\delta_{K'}(i,c)=\delta_K(i,c)+\sigma\bigl(1+\tfrac{1-2\cdot[i\in C_c]}{n_c}\bigr)$: the shift adds $\sigma(1+1/n_c)$ to the squared distance from $i$ to every other cluster $c$ and only $\sigma(1-1/n_a)$ to the distance to its own cluster $a$. Hence the Voronoi rule ``assign $i$ to $\arg\min_c\delta_{K'}(i,c)$'' favours staying, and as $\sigma\to\infty$ every partition becomes Voronoi-stable, while $E_\spec$ and the single-move dynamics are unchanged. \status{proved here (two lines); explicit form of an effect noted by \citet{dhillon2005tr,dhillon2007}}
\end{proposition}
\begin{proof}
$(\sigma I\one_c)_i=\sigma[i\in C_c]$ and $\one_c^{\!\top}\sigma I\one_c=\sigma n_c$; substitute in $\delta$.
\end{proof}
\noindent \citet{dhillon2005tr,dhillon2007} already discuss how the diagonal shift affects the practical performance of kernel $k$-means; the proposition is the explicit mechanism and E3 quantifies it. So ``kernel $k$-means with Lloyd's iteration'' is not one algorithm for a spring energy but one per PSD representative; the single-move relaxation is the representative-free one. For springs of negative type the natural representative is $\widetilde K$ of Theorem~\ref{thm:main}(f), which for Hooke springs is $2\langle x,y\rangle$ and reproduces ordinary Lloyd $k$-means; the experiments use it, and use the naively shifted $-A+\sigma I$ only as a control (E3).

\section{What leaves the family}\label{sec:escape}

Theorem~\ref{thm:main} covers every objective that is a sum over within-cluster \emph{pairs} of a fixed function of the data, weighted by a function of the cluster size in $\{1,1/n_c\}$. An ingredient can escape only by breaking one of these three features: the sum over pairs (many-body terms), the fixedness of the pair function (parameters that depend on the assignment), or the normalisation. A fourth route replaces the minimisation over partitions by a dynamics on positions. For the first three, specific instances can be certified exactly.

\begin{theorem}[Exact separation on explicit configurations; computer-assisted]\label{thm:e4}
Let $P_1$ and $P_2$ be the configurations of $\EfourN$ points in $\Q^2$
\begin{gather*}
P_1=\{\EfourPoints\},\\
P_2=\{\EfourPointsB\},
\end{gather*}
drawn once from a dedicated generator (seed \EfourSeed), and let $k=\EfourK$, so that each has $\EfourPartitions$ partitions. Call a partition function $F$ \emph{$w$-pairwise} if $F(C)=\sum_cw(n_c)\sum_{i<j\in C_c}A_{ij}+a$ for some symmetric $A$ and constant $a$ ($\EfourUnknowns$ unknowns). Then, with squared distances $s_{ij}=\lVert x_i-x_j\rVert^2$, the following holds on both configurations (in parentheses the relative residual $\lVert b-\Pi b\rVert/\lVert b\rVert$ on $P_1$; $P_2$, where $b\in\Q^{\EfourPartitions}$ lists the values of $F$ and $\Pi$ is the orthogonal projection onto the $w$-pairwise functions):
\begin{center}\footnotesize
\setlength{\tabcolsep}{4pt}
\begin{tabular}{lcc}
\toprule
partition function & $w_\spec$-pairwise? & $w_\tot$-pairwise?\\
\midrule
$\sum_c \frac1{n_c}\sum d_{ij}^2$ (Hooke, per particle) & yes (0) & no (0.023) \\
$\sum_c \sum d_{ij}^2$ (Hooke, total) & no (0.180) & yes (0) \\
$\sum_c \frac1{n_c}\sum (d_{ij}^2-\rho)^2$, fixed $\rho$ & yes (0) & no (0.042) \\
$\sum_c \frac1{n_c}\sum (d_{ij}^2-\rho_c)^2$, $\rho_c$ = cluster mean & no (0.181) & no (0.110) \\
$\sum_c \frac1{n_c}\sum_{i<j<l} 4\,\mathrm{area}(ijl)^2$ (three-body) & no (0.327) & no (0.056) \\
$\sum_c \binom{n_c}{2}^{-1}\sum d_{ij}^2$ (per spring) & no (0.203) & no (0.148)
\\
\bottomrule
\end{tabular}
\end{center}
where $\rho$ is the mean of all $s_{ij}$, $\rho_c$ the mean of $s_{ij}$ over the pairs of cluster $c$, and $\mathrm{area}(ijl)$ the area of the triangle $x_ix_jx_l$. Each row is one function $F$, tested against both families.
\end{theorem}
\begin{proof}
For fixed $w$ the $w$-pairwise functions form a linear subspace of $\R^{\EfourPartitions}$ spanned by the columns of an explicit rational design matrix (rank $\EfourRankSpecific$ for $w_\spec$, $\EfourRankTotal$ for $w_\tot$, on both configurations). Membership is decided by exact least squares in rational arithmetic (\texttt{fractions.Fraction}): a zero residual exhibits $A$ and $a$; a non-zero residual is a certificate that $F$ is not in the span. The script \path{experiments/identity_check.py} (E4) performs the computation; the relative residuals are reported in the table.
\end{proof}
\noindent The two positive controls (Hooke and fixed rest length under their own normalisation, total Hooke under $w_\tot$) confirm that the test is not blind, and the yes/no pattern agrees on the two configurations in \EfourCellsSame{} of the \EfourCells{} cells (a sanity check: one configuration suffices logically). The table says that the two normalisations are different objectives (each fails to be pairwise under the other's weight), and that the three functionals tested, all with the per-particle weight $1/n_c$ inside $F$ (rows 4--6), are neither $w_\spec$- nor $w_\tot$-pairwise; their total-weight versions were not tested. The per-spring average \emph{is} a pairwise sum with the size weight $\binom{n_c}{2}^{-1}$; the certificate only excludes the weights $1$ and $1/n_c$. By Cauchy--Binet applied to the rows $(1,x_i^{\!\top})$, $\sum_{i<j<l\in C}4\,\mathrm{area}(ijl)^2=|C|\det S_C$ with $S_C$ the scatter matrix of $C$, so the three-body row is $\sum_c\det S_c$, a quartic in the data. The theorem certifies these three functionals, not the classes ``many-body terms'' or ``assignment-dependent rest lengths'', and the class statements would be false: restricted to partitions whose clusters all have at least three points, $\sum_cw(n_c)/(n_c-2)\sum_{i<j<l\in C_c}(s_{ij}+s_{jl}+s_{il})=\sum_cw(n_c)\sum_{i<j\in C_c}s_{ij}$.

\emph{Values, not minimisers.} The certificate concerns objective values up to an additive constant. It does not show that the minimisers differ from those of every pairwise objective: with $A$ free, every partition is the unique minimiser of some $w_\tot$-pairwise objective ($A_{ij}=-1$ within its blocks, $0$ otherwise), and a strictly increasing non-affine transform of a pairwise objective in general fails the test while keeping its minimisers. Whether these ingredients can produce an \emph{observable} mechanical effect is therefore a separate question that the theorem does not answer.

\paragraph{The dynamic route.} Several physics-inspired algorithms do not minimise an energy over partitions but move the \emph{positions}: gravitational clustering~\citep{wright1977}, mean shift~\citep{fukunaga1975,comaniciu2002} (query points ascend a fixed kernel density estimate) and its blurring variant, in which all data points move~\citep{cheng1995}, quantum clustering~\citep{horn2002}, and the finite-temperature correlations of a Potts model~\citep{blatt1996}. Their output is a function of the trajectory, and Theorem~\ref{thm:main} says nothing about it.

\begin{problem}[Dynamic mechanical clustering]\label{prob:dyn}
Let $V\colon(0,\infty)\to\R$ be a pair potential, let $z(t)$ solve $\dot z_i=-\sum_{j\ne i}V'(\lVert z_i-z_j\rVert)\,\frac{z_i-z_j}{\lVert z_i-z_j\rVert}$ with $z(0)=X$, and let $P_{T,\varepsilon}(X)$ be the single-linkage partition of $z(T)$ at threshold $\varepsilon$. Characterise the triples $(V,T,\varepsilon)$ for which there exist a fixed function $a\colon[0,\infty)\to\R$ and a normalisation $w$ such that, for every $X$, $P_{T,\varepsilon}(X)$ minimises $\sum_cw(n_c)\sum_{i<j\in C_c}a(\lVert x_i-x_j\rVert)$ over partitions into the number of clusters it produces.
\end{problem}

\noindent For the Gaussian attraction $V(d)=-e^{-d^2/2h^2}$ the flow is a continuous-time, unnormalised form of Gaussian \emph{blurring} mean shift~\citep{cheng1995}, not of standard mean shift. We expect no such pair $(a,w)$ to exist for generic $(T,\varepsilon)$, but have not tested it. Whatever the answer, a mechanical proposal of this kind has to be compared with (blurring) mean shift, not with $k$-means, and the comparison must involve the trajectory, not the energy minimum.

\section{Computational confirmation}\label{sec:experiments}

\paragraph{Data and potentials.} Five datasets: blobs, two moons and two circles from scikit-learn's generators (\Nsynth{} points each in $\R^2$; \texttt{make\_blobs} with three centres and std $0.9$, $k=3$; \texttt{make\_moons} with noise $0.06$ and \texttt{make\_circles} with factor $0.45$ and noise $0.05$, $k=2$), a random subset of $\Ndigits$ digits from \texttt{load\_digits} (64 features scaled to $[0,1]$, $k=10$) and the standardised iris data (\Niris{} points, 4 features, $k=3$). The six potentials of Section~\ref{sec:family}, with scales set from the data before any run: $\ell$ is the median distance to the seventh nearest neighbour, $r$ the tenth percentile of all pairwise distances, $s$ the median of $\lVert x\rVert^2$. In all, $\NumDatasets\times\NumPotentials=\EoneConfigs$ configurations. Seeds: \MetaSeed{} (E1, E2), \MetaSeedRelax{} (E3) and \EfourSeed{} (E4).

\paragraph{E1: identical objective values.} \emph{Criterion (predefined):} for every tested partition and every potential, $|E_\spec-(\tfrac12J_{K'}+\mathrm{const})|\le10^{-9}\max(1,|E_\spec|)$ with the constant of Theorem~\ref{thm:main}, for the three representatives $K'\in\{-A,\;-A+\sigma I,\;\widetilde K\}$; pass means $100\%$. The partitions are \EoneRandomPartitions{} uniformly random ones per configuration, the outputs of Lloyd and of the single-move relaxation on $\widetilde K$ (minimally shifted when indefinite), and the ground-truth labels. \emph{Result:} \EonePassed/\EoneChecks{} checks passed, worst relative error $\EoneMaxRelErr$ (the shifted representative uses $\sigma$ up to $\EoneMaxSigma$, so its identity involves cancellation); the Hooke energy equals the $k$-means SSE to $\EoneHookeSSEErr$. On equal-size partitions $E_\tot=(n/k)E_\spec$ to $\EoneEqualSizeErr$, while on random partitions the ratio $E_\tot/E_\spec$ ranges over $[\EoneRatioMin,\EoneRatioMax]$: the normalisation is a genuine choice. The spectrum of $\widetilde K$ (per configuration in \path{results/tables_identity.md}): for the \EoneCndConfigs{} configurations with a potential of negative type it is PSD in \EoneCndPsd{} (numerically, $\lambda_{\min}/\lambda_{\max}\ge-\EoneCndWorstRatio$), and for the rest-length potential it is indefinite in \EoneRestIndefinite/\EoneRestConfigs.

\paragraph{E2: same optimiser, same initialisation, same trajectory.} \emph{Criterion:} a brute-force spring relaxation that recomputes the total energy from scratch for every candidate move (it sees only springs) and the kernel implementation of Hartigan's method via~\eqref{eq:delta} (it sees only $K$), started from the same random labels with the same sweep order, must produce identical label vectors after every sweep and final energies equal to $10^{-9}$ relative, in all runs. Subsets of $\EtwoNsub$ points, three initialisations per configuration. \emph{Result:} \EtwoIdentical/\EtwoRuns{} identical trajectories (mean \EtwoMeanSweeps{} sweeps, at most \EtwoMaxSweeps), worst final-energy discrepancy $\EtwoMaxRelErr$. In addition, hub-anchored Hooke springs equilibrated alternately (Remark~\ref{rem:lift}(b)) reproduced scikit-learn's Lloyd $k$-means from the same initial centroids in \EtwoAnchoredIdentical/\EtwoAnchoredRuns{} runs, with inertia equal to the spring energy to $\EtwoAnchoredMaxErr$.

\paragraph{E3: relaxation versus kernel $k$-means on the same energy.} Five optimisers see the same kernel: \emph{lloyd} (batch Lloyd on $\widetilde K$, minimally shifted only when indefinite, \EthreeShiftedConfigs{} configurations), \emph{lloyd+relax} (Lloyd, then single-move polish), \emph{relax} (single-particle relaxation from the initial labels), \emph{anneal} (Metropolis on the spring energy, \EthreeAnnealSweeps{} sweeps of geometric cooling from the median $|\Delta E|$ of random moves down by $10^{-3}$, then relaxation~\citep{kirkpatrick1983}; \citet{hofmann1997} anneal the same energy deterministically) and the control \emph{lloyd (naive)} (Lloyd on $-A+\sigma I$, Proposition~\ref{prop:shift}). Initialisations are shared: $\EthreeR$ random labelings and $\EthreeR$ kernel $k$-means++ seedings~\citep{arthur2007} per configuration; annealing uses the first $\EthreeRanneal$ random ones. $E^*$ is the best energy found by any method in a configuration. \emph{Criteria (predefined):} (C3a) every relaxation and annealing fixed point is Voronoi-stable for the kernel used (Proposition~\ref{prop:hartigan} predicts $100\%$); (C3b) the relaxation is called distinguishable from the kernel family only if its best-of-restarts energy differs from that of lloyd+relax by more than $10^{-6}$ relative, either sign being reported since it would be a statement about optimisers of one function; (C3c) the fraction of Lloyd fixed points that are not single-move stable is reported. \emph{Disclosure:} a pilot run used $-A+\sigma I$ as the only Lloyd kernel and found it inert (Proposition~\ref{prop:shift}); $\widetilde K$ then replaced it as the baseline, which favours the kernel side of the comparison, and the naive shift was kept as a control. C3a's wording was updated to name the kernel used; its prediction ($100\%$, valid for every PSD representative by Proposition~\ref{prop:hartigan}(ii)), C3b, C3c and all thresholds were written before the pilot and not changed; every number reported comes from the final code. \emph{Results} (Table~\ref{tab:e3}; per configuration in \path{results/tables_relaxation.md}, and the relative excess energies of the four Lloyd and relaxation optimisers per dataset and potential in \path{figures/energies.pdf}): all \EthreeVoronoiRelax/\EthreeRunsRelax{} relaxation and \EthreeVoronoiAnneal/\EthreeRunsAnneal{} annealing fixed points are Voronoi-stable (C3a). Under C3b, \EthreeDistinguishable{} of \EthreeConfigs{} configurations are distinguishable (\EthreeRelaxBetter{} in favour of the relaxation, \EthreeLloydRelaxBetter{} against; largest relative difference $\EthreeMaxAbsDiff$; the split is optimiser noise: the \texttt{--fast} run gives 3--3 with the same 6/30); whenever the two best energies agree the best partitions coincide (ARI $=1$; \EthreeAgreeButDifferent{} exceptions), and overall ARI $=1$ in \EthreeAriOne/\EthreeConfigs. The relaxation reaches $E^*$ in \EthreeReachRelax{} configurations, lloyd+relax in \EthreeReachLloydRelax, Lloyd alone in \EthreeReachLloyd{} and the naive shift in \EthreeReachNaive{} (annealing, with \EthreeRanneal{} restarts only, is not comparable on this count: \EthreeReachAnneal). Only \EthreeLloydHartiganPct\% of Lloyd fixed points on the canonical kernel are single-move stable, \EthreeNaiveHartiganPct\% on the naive shift (C3c); median excesses and the measured single-move stability of the single-move methods are in \path{results/tables_relaxation.md}. The best energies were recomputed spring by spring, worst discrepancy $\EthreeDirectCheck$. E3 is descriptive (counts over \EthreeConfigs{} configurations, modest restart budgets); no inferential claim is made.

\begin{table}[t]
\centering\scriptsize
\setlength{\tabcolsep}{3pt}
\begin{tabular}{lcccccccccc}
\toprule
 & \multicolumn{4}{c}{configurations reaching $E^*$ (of 5)} & \multicolumn{3}{c}{restarts reaching $E^*$ (\%)} & Lloyd f.p. & best\\
\cmidrule(lr){2-5}\cmidrule(lr){6-8}
potential & relax & lloyd+relax & lloyd & naive & relax & lloyd+relax & lloyd & 1-move stable & ARI$=1$\\
\midrule
$d^2$ (Hooke) & 4/5 & 5/5 & 3/5 & 1/5 & 54 & 57 & 31 & 46/100 & 4/5 \\
$d^2$ on the lift $(x,|x|^2/s)$ & 5/5 & 5/5 & 4/5 & 2/5 & 80 & 74 & 69 & 73/100 & 5/5 \\
$(d-r)^2$ (rest length) & 5/5 & 4/5 & 1/5 & 1/5 & 53 & 49 & 4 & 4/100 & 4/5 \\
$d$ (constant tension) & 5/5 & 5/5 & 2/5 & 1/5 & 42 & 43 & 30 & 43/100 & 5/5 \\
$1-e^{-d^2/2\ell^2}$ (saturating) & 4/5 & 2/5 & 2/5 & 0/5 & 4 & 2 & 2 & 12/100 & 2/5 \\
$\log(1+d^2/\ell^2)$ (soft) & 5/5 & 4/5 & 2/5 & 1/5 & 37 & 34 & 17 & 26/100 & 4/5
\\
\bottomrule
\end{tabular}
\caption{E3 by potential (five datasets each). ``Reaching $E^*$'' means within $10^{-6}$ relative of the best energy found by any method. ``Lloyd f.p.\ 1-move stable'': Lloyd fixed points that no single move improves (all relaxation fixed points are Lloyd fixed points, Proposition~\ref{prop:hartigan}(ii); 100/100 per potential). ``best ARI$=1$'': configurations in which the best partitions of relax and lloyd+relax coincide.}
\label{tab:e3}
\end{table}

\paragraph{Reading of E1--E3 for the research line.} If the screened formulations were of the pairwise kind of Definition~\ref{def:family}, the screening could not have found a mechanical effect: the spring energies \emph{are} the kernel objectives (E1), the mechanical implementation of the relaxation \emph{is} a kernel implementation (E2), and every relaxed state \emph{is} a kernel $k$-means fixed point (E3). Differences between optimisers of the same function are visible (Lloyd alone reaches $E^*$ less often than the single-move methods, consistent with \citet{franca2021}, who found Hartigan's method preferable to Lloyd's for kernel $k$-means; the naive shift is worse still), and they are the only differences there are.

\section{What can be claimed}\label{sec:claims}

Table~\ref{tab:claims} collects the mathematical claims with their status; it separates what is known, what is proved here and what was only measured (Section~\ref{sec:experiments}; ``verified'' there means reproduced by the scripts of Appendix~\ref{app:repro} with fixed seeds, not proved). Most of what closes the screening is literature.

\begin{table}[t]
\centering\scriptsize
\setlength{\tabcolsep}{3pt}\setlength{\extrarowheight}{0.5pt}
\begin{tabular}{>{\raggedright\arraybackslash}p{0.56\linewidth}>{\raggedright\arraybackslash}p{0.41\linewidth}}
\toprule
claim & status\\
\midrule
$J$ is invariant under diagonal shift (up to $\sigma(n-k)$), constant and separable terms (Lemma~\ref{lem:inv}) & elementary; cf.\ \citealp{roth2003,dhillon2007}\\
Per-particle spring energy $=$ pairwise-clustering cost $=\tfrac12J_{-A}+\tfrac12(n-k)\phi(0)$, PSD after a shift (Thm.~\ref{thm:main}(b)) & literature (\citealp{hofmann1997}; constant-shift embedding \citealp{roth2003}; ratio association $\leftrightarrow$ kernel $k$-means \citealp{dhillon2005tr,dhillon2007}); constant made explicit\\
Zero-rest-length Hooke energy $=$ $k$-means SSE (Thm.~\ref{thm:main}(c)) & classical\\
Total energy / released tension $=$ min-sum $k$-clustering / max-$k$-cut (Thm.~\ref{thm:main}(d)) & elementary\\
Every kernel $k$-means objective is a Hooke energy of a lifted configuration (Thm.~\ref{thm:main}(e)) & classical\\
Negative-type springs: dataset-independent PD kernel; Hooke springs after a lift and a $1/\sqrt2$ rescaling (Thm.~\ref{thm:main}(f)) & literature (\citealp{bcr1984,sejdinovic2013,franca2021}); written for spring potentials\\
$d^2$ is of negative type; Bernstein transforms preserve negative type (Prop.~\ref{prop:bernstein}(i)--(iii)) & classical / literature\\
$(d-r)^2$ is not of negative type; no PD kernel $\widetilde K+f(x)+f(y)+a+\sigma\delta_{xy}$ (Prop.~\ref{prop:bernstein}) & proved here (elementary); other representatives open\\
Lifting only changes the kernel; anchored Hooke springs are Lloyd $k$-means (Remark~\ref{rem:lift}) & trivial / classical\\
Single-particle relaxation is Hartigan's method in kernel space (Prop.~\ref{prop:hartigan}(i)) & literature (kernel $k$-groups, \citealp{franca2021})\\
Its fixed points are Voronoi-stable for every PSD representative (Prop.~\ref{prop:hartigan}(ii)) & proved here (kernel form of \citealp{telgarsky2010}; priority not checked)\\
Lloyd's rule depends on the PSD representative; large shifts freeze it (Prop.~\ref{prop:shift}) & proved here; effect noted by \citet{dhillon2005tr,dhillon2007}\\
The per-particle triangle-area three-body term and cluster-mean rest length, and the per-spring average, are not $w$-pairwise for $w\in\{1,1/n_c\}$; the two normalisations differ (Thm.~\ref{thm:e4}; values, not minimisers) & proved here (exact arithmetic, two explicit configurations)\\
Dynamic (position-relaxation) mechanical clustering (Problem~\ref{prob:dyn}) & open; not tested\\
Any statement about clustering quality on real data & not addressed\\
\bottomrule
\end{tabular}
\caption{Mathematical claims and their status.}
\label{tab:claims}
\end{table}

\paragraph{Limitations.} The family of Section~\ref{sec:family} is a reconstruction from the line's public summary; if the original screening used many-body terms, assignment-dependent parameters or a dynamics on positions, Section~\ref{sec:escape} applies and the theorem does not. The theorem is about objectives; it does not say that all optimisers of a spring energy behave alike, and E3 shows that they do not. The normalisation family is $\{1,1/n_c\}$. Theorem~\ref{thm:e4} certifies three functionals on two configurations and concerns objective values, not minimisers; it says nothing about classes of terms or about how far from pairwise those energies are on typical data. The Bernstein composition result is cited, not re-proved in full. The attributions to \citet{franca2021} and to the diagonal-shift discussion of \citet{dhillon2005tr,dhillon2007} were checked against abstracts and bibliographic indices, not against the full texts. E3 uses one implementation of each optimiser with modest restart budgets; its numbers describe optimiser luck on five datasets, not a benchmark. Nothing is claimed about the dynamic route beyond its formulation.

\paragraph{Next steps for the research line.} (a) Confirm with the author of the line that the screened formulations were of the pairwise form of Definition~\ref{def:family}; if not, the emphasis moves to Section~\ref{sec:escape}. (b) Settle Open problem~\ref{prob:dyn} for the Gaussian flow, by configurations on which no fixed pair function reproduces $P_{T,\varepsilon}$ as a minimiser (Theorem~\ref{thm:e4} tests values, so it does not apply directly), or by a representation. (c) Characterise the size weights $w$ for which $\sum_cw(n_c)\sum_{i<j\in C_c}A_{ij}$ is a weighted kernel $k$-means objective~\citep{dhillon2007}. (d) Decide whether Lemma~\ref{lem:inv} lists all invariances of $J$ (this settles the rest-length representative). (e) If (a) is confirmed, the line's summary can be updated from ``no mechanical component survived'' to ``no mechanical component of the pairwise kind can survive'', with the open ingredients named.

\appendix
\section{Reproducibility}\label{app:repro}
{\sloppy \path{experiments/mechanical.py} implements the potentials, energies, kernel forms and optimisers; \path{experiments/identity_check.py} runs E1, E2 and E4 and \path{experiments/relaxation.py} runs E3 (both take \texttt{--fast}); \path{experiments/make_numbers.py} converts \path{results/*.json} into the macros and table bodies used here, so every number in the text is generated, and \path{results/*.md} hold the full tables. Seeds \MetaSeed{} (E1, E2), \MetaSeedRelax{} (E3) and \EfourSeed{} (E4, own generator); Python \MetaPython, NumPy \MetaNumpy, SciPy \MetaScipy~\citep{virtanen2020}, scikit-learn \MetaSklearn~\citep{pedregosa2011}; \IdentitySeconds{} s (E1, E2, E4) and \RelaxSeconds{} s (E3) on shared cores. Implementation details (move thresholds, tolerances, empty clusters, annealing proposals, the E4 generator and the exact elimination over $\Q$) are listed in the README.\par}

\bibliographystyle{abbrvnat}
\section*{References}
\renewcommand{\bibsection}{}
\raggedcolumns
\begin{multicols}{2}
\bibliography{refs}
\end{multicols}
\end{document}
